The main reproduction run constructed 37 positive certificates and
translated each to the matched nested checker and the independently
retrieved existing four-family checker. Each ran in all three programs
under both ordinary Python and \code{python -O}. An additional 22 negative
inputs ran in both modes against the prefix checker. This is 266 subprocess
expectations, with zero mismatches. The positive cases include zero and
successor computation, derived symmetry, nested paired binders, and a
nontrivial equation-induction proof of
\[
 \I(y,0;u.Su)=y
\]
followed by specialization. The negative cases include incorrect induction
steps, escaped eigenparameters, malformed or out-of-scope terms, invalid
clause indices, nonintegral binder counts, and compound premises used as
equations.

There were also 2001 checks that opening commutes with expansion to the
nested representation. A separately written audit used environment-based
numerical semantics, rather than another syntactic substitution routine,
for 400 matrix openings, 400 witness openings, and 300 iterator-body
substitutions. It passed 20 additional rejection cases and 5 accepted
examples. The scope attacks include a forbidden parameter beneath a dead
zero-iterator body and beneath a formula binder; dead code is still
checked for scope.

The deterministic generators, concrete certificates, full results, and
source hashes are included. The independent semantic checks use bounded
samples and a finite test interpretation for formula-opening comparisons;
they are not asserted to be a model of the entire arithmetic theory.
Written derivation translations supply the metatheorems, while these
checks supply evidence that the implementations follow them.

From the package directory, run
\begin{lstlisting}
python3 code/reproduce_checker.py
python3 code/audit_flat.py code/flat_control.py
python3 code/flat_control.py < code/examples/equation_induction.json
\end{lstlisting}
The last command checks the supplied induction certificate. The example
filenames and complete format description are also listed in
\code{README.txt} and \code{code/FORMAT.txt}.
